\section*{Phase I SARS-CoV-2 HRSM Workplan}

\textbf{Generated:} 2026-06-03

\subsection*{Purpose}

Phase I tests whether persistent SARS-CoV-2 infection can be represented as a viral HRSM state space defined by adaptive reserve, rebound capacity, swarm coherence, and retained evolutionary memory.

\subsection*{Model comparison}

The first model comparison is deliberately narrow:

\begin{itemize}
\item Markovian current-state model: \(H_v, R_v, S_v\).
\item Non-Markovian memory model: \(H_v, R_v, S_v, M_v\).
\end{itemize}

\subsection*{Ranked candidate datasets}

\subsubsection*{1. Kemp et al. chronic infection under convalescent plasma}
\begin{itemize}
\item \textbf{Paper:} Kemp et al., Nature 2021
\item \textbf{URL:} https://www.nature.com/articles/s41586-021-03291-y
\item \textbf{Data anchor:} Check paper Data Availability and supplementary materials
\item \textbf{Reported timepoints:} 23
\item \textbf{Reported duration days:} 101
\item \textbf{Data strength:} very\_high
\item \textbf{Recommended role:} primary\_phase1\_test\_case
\item \textbf{Audit note:} Best first target. Whole-genome ultra-deep sequencing across 23 timepoints. Strong treatment-pressure structure.
\end{itemize}

\subsubsection*{2. Schmidt et al. 521-day persistent infection}
\begin{itemize}
\item \textbf{Paper:} Schmidt et al., npj Genomic Medicine 2025
\item \textbf{URL:} https://www.nature.com/articles/s41525-025-00463-x
\item \textbf{Data anchor:} ENA PRJEB77414
\item \textbf{Reported timepoints:} 5
\item \textbf{Reported duration days:} 521
\item \textbf{Data strength:} high
\item \textbf{Recommended role:} long\_memory\_case
\item \textbf{Audit note:} Excellent memory branch because duration is extreme. Fewer timepoints than Kemp, but long temporal span is useful for M\_v.
\end{itemize}

\subsubsection*{3. Weigang et al. immunosuppressed within-host evolution}
\begin{itemize}
\item \textbf{Paper:} Weigang et al., Nature Communications 2021
\item \textbf{URL:} https://www.nature.com/articles/s41467-021-26602-3
\item \textbf{Data anchor:} Check Data Availability, supplementary files, and sequence repositories
\item \textbf{Reported timepoints:} to\_audit
\item \textbf{Reported duration days:} to\_audit
\item \textbf{Data strength:} medium\_high
\item \textbf{Recommended role:} immune\_escape\_validation\_case
\item \textbf{Audit note:} Good immune-escape trajectory. Needs accession-level check before becoming a primary case.
\end{itemize}

\subsubsection*{4. Harari et al. chronic infection adaptive evolution drivers}
\begin{itemize}
\item \textbf{Paper:} Harari et al., Nature Medicine 2022
\item \textbf{URL:} https://www.nature.com/articles/s41591-022-01882-4
\item \textbf{Data anchor:} Check supplementary tables and public sequence identifiers
\item \textbf{Reported timepoints:} multi\_case
\item \textbf{Reported duration days:} multi\_case
\item \textbf{Data strength:} medium
\item \textbf{Recommended role:} cross\_case\_context\_and\_robustness
\item \textbf{Audit note:} Important chronic-infection context across 27 infections, but may be less clean as the first mechanistic HRSM case.
\end{itemize}

\subsubsection*{5. Ghafari et al. community surveillance persistent infections}
\begin{itemize}
\item \textbf{Paper:} Ghafari et al., Nature 2024
\item \textbf{URL:} https://www.nature.com/articles/s41586-024-07029-4
\item \textbf{Data anchor:} Check study data access conditions and supplementary metadata
\item \textbf{Reported timepoints:} large\_surveillance
\item \textbf{Reported duration days:} 30\_to\_60\_plus
\item \textbf{Data strength:} medium
\item \textbf{Recommended role:} background\_population\_reference
\item \textbf{Audit note:} Excellent for population-scale framing, not the first mechanistic within-host HRSM case.
\end{itemize}

\subsection*{Phase I execution sequence}
\begin{enumerate}
\item Verify accession and supplementary data availability for Kemp et al.
\item Build a case-level sample table for the Kemp trajectory.
\item Extract sequence, timepoint, treatment, and abundance proxies.
\item Compute first-pass \(H_v\), \(R_v\), \(S_v\), and \(M_v\) axes.
\item Run leakage audit, axis correlation audit, and sampling-duration audit.
\item Compare current-state and memory-augmented models.
\item Repeat only after the first case pipeline is stable.
\end{enumerate}

\subsection*{Failure rules}
\begin{itemize}
\item Do not treat dataset priority as proof of data usability.
\item Do not use treatment status as both an axis component and a predicted endpoint.
\item Do not let \(M_v\) collapse into ordinary genetic divergence.
\item Do not interpret viral persistence as organism-level homeostasis.
\end{itemize}